\documentclass[12pt]{article}

% -------------------------------------------------------
% Packages
% -------------------------------------------------------
\usepackage[margin=1in]{geometry}
\usepackage[T1]{fontenc}
\usepackage{lmodern}
\usepackage{amsmath,amssymb,amsthm}
\usepackage{booktabs}
\usepackage{hyperref}
\usepackage{xcolor}
\usepackage{listings}
\usepackage{graphicx}
\usepackage{microtype}
\usepackage{authblk}
\usepackage{caption}
\usepackage{float}
\usepackage{enumitem}
\usepackage{array}
\usepackage{multirow}

% -------------------------------------------------------
% Theorem environments
% -------------------------------------------------------
\newtheorem{theorem}{Theorem}
\newtheorem{lemma}[theorem]{Lemma}
\newtheorem{proposition}[theorem]{Proposition}
\newtheorem{corollary}[theorem]{Corollary}
\newtheorem{observation}[theorem]{Observation}
\theoremstyle{definition}
\newtheorem{definition}[theorem]{Definition}
\newtheorem{remark}[theorem]{Remark}

% -------------------------------------------------------
% Code listings
% -------------------------------------------------------
\lstset{
  basicstyle=\ttfamily\small,
  frame=single,
  breaklines=true,
  numbers=left,
  numberstyle=\tiny,
  keywordstyle=\bfseries\color{blue!70!black},
  commentstyle=\color{green!50!black},
}

% -------------------------------------------------------
% Macros
% -------------------------------------------------------
\newcommand{\Gn}{G_n}
\newcommand{\pw}{\mathrm{pw}}
\newcommand{\tw}{\mathrm{tw}}
\newcommand{\Bell}{\mathrm{Bell}}
\newcommand{\maxfw}{\text{max\_fw}}
\newcommand{\nb}{n_{\mathrm{back}}}
\newcommand{\fs}{f_s}
\newcommand{\nxt}{\texttt{nxt}}
\newcommand{\curr}{\texttt{curr}}
\newcommand{\proxy}{\hat{s}}   % geometric state-count proxy

% -------------------------------------------------------
\begin{document}

\title{\textbf{Exact Counting of Hamiltonian Paths in Square-Sum Graphs:\\
An Iterative Engineering Approach}}

\author{Victor S. Miller}

\date{\today}

\maketitle

\begin{abstract}
For each positive integer~$n$, define the \emph{square-sum graph}~$\Gn$ to be
the undirected graph on vertex set $\{1,\ldots,n\}$ in which vertices $i$ and~$j$
are adjacent if and only if $i+j$ is a perfect square.
We describe an algorithmic journey to compute the exact number of Hamiltonian
paths in~$\Gn$ for all~$n \le 75$,
a computation previously limited to $n \le 56$ by memory constraints.

Starting from a \texttt{graphillion}-based ZDD enumeration paired with a
MaxSAT pathwidth ordering~\cite{Coudert}, we developed a compact
C~frontier dynamic-programming (DP) implementation.
A series of targeted optimisations — packed \texttt{uint64\_t} state
keys, dynamic table sizing, multi-edge batching, fused edge-and-elimination
sweeps, and a secondary vertex-ordering refinement via simulated annealing
on an empirically motivated cost function — reduced the per-$n$ DP time for
$n = 57$ from roughly $134$~seconds to around $30$~seconds while
simultaneously reducing peak memory from near the $16$\,GB machine limit
to a fraction of that.
Replacing the hash table by an extendible hash table with a radix-sorted
insertion buffer, and returning freed memory to the operating system, then
allowed $n=61$ in $97$ minutes in under $2$\,GB on a laptop.
Finally, a parallel sort--merge formulation of the DP step, in which
independent contributions are generated by worker threads, sorted by a hash
of the key, and merged by threads owning disjoint ranges of hash space,
carried the computation to $n=75$ on a $768$\,GB, $64$-core machine.
Storing every table in hash order balances the work without any model of
its cost.
The orderings now come from a solver based on Kobayashi's dynamic program for
pathwidth, which finds the exact pathwidth in seconds and represents a very
large family of pathwidth-optimal orderings as a DAG, over which a secondary
cost is minimised exactly.
\end{abstract}

% -------------------------------------------------------
\section{Introduction}
% -------------------------------------------------------

The \emph{square-sum graph}~$\Gn$ is a natural combinatorial object studied
in recreational and competitive mathematics (it is related to the question of
whether $\{1,\ldots,n\}$ can be arranged in a cyclic sequence in which every
adjacent pair sums to a perfect square).
Counting its Hamiltonian paths exactly is computationally hard in general — the
problem is \#P-complete — but for specific sparse graph families the structure
of the graph can be exploited.

Our starting point was an existing implementation that used the
\texttt{graphillion} library~\cite{graphillion} (which represents sets of
graphs compactly as Zero-suppressed Binary Decision Diagrams, or ZDDs) together
with a MaxSAT-based exact pathwidth ordering (described in
Section~\ref{sec:ordering}).
This pipeline could handle $n \le 56$ before exhausting the available $16$\,GB
of RAM.

The rest of this paper is organised as follows.
Section~\ref{sec:background} establishes background on frontier DPs and
pathwidth.
Section~\ref{sec:induction} addresses a proposed inductive counting approach
and shows why the frontier DP subsumes it.
Sections~\ref{sec:dp}--\ref{sec:memory} describe the sequence of
engineering improvements, each motivated by empirical profiling.
Section~\ref{sec:eh} describes the extendible-hashing version used up to
$n=61$ on a laptop, and Section~\ref{sec:sm} the parallel sort--merge version
used for $62\le n\le 75$.
Section~\ref{sec:kpw} describes the Kobayashi-based ordering solver that
replaced the MaxSAT ordering and the annealing.
Section~\ref{sec:results} summarises the results obtained, and
Section~\ref{sec:conclusion} concludes.

% -------------------------------------------------------
\section{Background}
\label{sec:background}
% -------------------------------------------------------

\subsection{The square-sum graph}

$\Gn$ has $n$ vertices and, for the sizes we study, maximum degree around~$6$.
It is always connected for $n \ge 15$ and remains sparse: $|E(G_{57})| = 121$.
Selected values of $|\text{Ham}(G_n)|$ (undirected Hamiltonian paths) include:

\begin{center}
\begin{tabular}{rrr}
\toprule
$n$ & $|\text{Ham}(G_n)|$ & DP time (final implementation)\\
\midrule
15 & 1 & $<1$\,ms\\
25 & 10 & $<1$\,ms\\
35 & 17{,}175 & $<50$\,ms\\
40 & 39{,}550 & $<50$\,ms\\
57 & 68{,}092{,}497{,}615 & $\approx 30$\,s\\
\bottomrule
\end{tabular}
\end{center}

\subsection{Pathwidth and the frontier DP}
\label{sec:frontier}

The \emph{pathwidth} $\pw(G)$ of a graph~$G$ is the minimum, over all linear
orderings of~$V(G)$, of the maximum number of vertices that are
\emph{active} (i.e.\ introduced but not yet eliminated) at any step of
that ordering.
Formally, given an ordering $\sigma = (v_1,\ldots,v_n)$ and a step $i$,
the \emph{frontier width} at step~$i$ is
\[
  f_i(\sigma) = \bigl|\{v_j : j \le i,\; \exists\, k > i \text{ s.t.\ }(v_j,v_k)\in E\}\bigr|,
\]
and $\pw(G) = \min_\sigma \max_i f_i(\sigma)$.

The frontier DP processes vertices in the order $\sigma$.
At each step it maintains a table of \emph{states} describing which vertices
are active (the \emph{frontier}), how the selected edges partition them into
path segments, and a count of how many orderings of edges lead to each partition.
The number of distinct canonical states at any step is at most $\Bell(w)$,
the $w$-th Bell number, where $w$ is the current frontier width.
Since Bell numbers grow super-exponentially, minimising the maximum frontier
width (i.e.\ finding the pathwidth ordering) is critical.

\subsection{MaxSAT pathwidth ordering}
\label{sec:ordering}

The pathwidth ordering is computed by encoding the vertex-separation problem as
a weighted MaxSAT instance following the formulation of
Coudert~\cite{Coudert} and solving with the RC2 solver from
\texttt{pysat}.
This gives the exact pathwidth and an optimal ordering, which is then used as
input to the frontier DP.

For the graphs we study, typical values are $\pw(G_{57}) = 11$ and
$\pw(G_{59}) = 13$.
The MaxSAT solver requires $20$--$500$~seconds depending on~$n$.

% -------------------------------------------------------
\section{Why an Inductive Approach Is Equivalent to the Frontier DP}
\label{sec:induction}
% -------------------------------------------------------

An intuitive attempt at a faster algorithm proceeds as follows.
Define $p(n,s,t)$ to be the number of directed Hamiltonian paths in~$\Gn$
from~$s$ to~$t$.
When forming $G_{n+1}$ from~$G_n$ by adding vertex~$n+1$, paths in
$G_{n+1}$ either have $n+1$ as an endpoint (handled by summing the
appropriate column/row of the $p$-matrix) or have $n+1$ in the interior
(requiring the path to be split at~$n+1$).

The interior case requires not $p(n,s,t)$ but the count of pairs of
vertex-disjoint subpaths covering all of $\{1,\ldots,n\}$, which is a
$4$-dimensional table.
Extending to $G_{n+2}$ then requires a $6$-dimensional table, and so on.
The induction does not close.

\begin{observation}
The frontier DP is exactly a systematic implementation of the above induction.
The DP state at any step encodes all possible $k$-path covers of the vertices
introduced so far (for any $k \ge 1$), canonically.
The minimum number of distinct covers at any step, over all vertex orderings,
is bounded by $\Bell(\pw(G))$.
\end{observation}

% -------------------------------------------------------
\section{The Frontier DP: Correctness and Implementation}
\label{sec:dp}
% -------------------------------------------------------

\subsection{State encoding}

Each state is a tuple $(\sigma_0,\ldots,\sigma_{w-1}, c)$ where:
\begin{itemize}[nosep]
  \item $w$ is the current frontier size;
  \item $\sigma_i \in \{-1, 0, 1, 2, \ldots\}$ encodes the connectivity of
        the $i$-th frontier vertex: $0$ = unvisited; $L > 0$ = endpoint of
        partial path-segment~$L$; $-1$ = interior (degree 2, cannot take more edges);
  \item $c \in \{0,1\}$ counts \emph{closed chains} — path segments both of
        whose endpoints have left the frontier.
\end{itemize}

States are canonicalised by relabelling positive integers in order of first
appearance, so that states differing only by relabelling collapse into one.

\subsection{A subtle correctness bug and its fix}

A na\"ive frontier DP canonicalises only the connectivity tuple, ignoring
closed chains.
This causes a silent correctness failure: when two \emph{half-external} chains
(each with one endpoint that has already left the frontier) are merged by adding
an edge, both labels disappear from the state tuple — indistinguishable from
the legitimate ``all interior'' state.
These phantom states accumulate and are counted as valid Hamiltonian paths.

For $n = 15$ the na\"ive DP returns $11$ instead of the correct $1$.

\textbf{Fix.} Track $c \in \{0,1\}$ as an extra field.
Increment $c$ whenever a chain closes (via either elimination or a half-external
merge); discard states with $c > 1$ immediately; discard states with $c = 1$
that close at an intermediate step.

% -------------------------------------------------------
\section{From Python to C/cffi: The v1 Implementation}
\label{sec:c}
% -------------------------------------------------------

The Python reference implementation was correct but orders of magnitude slower
than \texttt{graphillion}.
We implemented the frontier DP in C, compiled as a shared library loaded via
\texttt{cffi}.

The v1 implementation used:
\begin{itemize}[nosep]
  \item 12-byte array keys with FNV-1a hashing;
  \item malloc/calloc/free on \emph{every} edge transition and every elimination
        (approximately $200{+}$ allocations for $n = 57$);
  \item a \texttt{used\_list} for O(\textit{cnt}) clears.
\end{itemize}

For $n = 57$ with the optimal pathwidth ordering this gave a DP time of
approximately $134$~seconds.

% -------------------------------------------------------
\section{Packed Keys and Dynamic Table Sizing: v2}
\label{sec:v2}
% -------------------------------------------------------

\subsection{Packed \texttt{uint64\_t} state keys}

Since the pathwidth of the graphs we study is at most $\sim\!15$, we can
encode each state as a single 64-bit integer: four bits per frontier slot
(values $1 = $ interior, $2 = $ unvisited, $k+2 = $ label $k$), one bit for
$c$, and one marker bit.
This replaces 12-byte array operations with single-integer hash/compare, and
canonicalisation with a nibble-loop on a register value.

The encoding fits frontier sizes up to $15$ (4 $\times$ 15 + 2 flag bits = 62
bits $\le$ 64).
Larger pathwidths would require a wider key type.

\subsection{Dynamic table sizing}

The v1 code kept tables at peak capacity even after the frontier shrank,
causing sequential scans proportional to the \emph{historical peak} rather
than the current state count.
The replacement function \texttt{ht\_ensure\_clear} sizes the table to
exactly $\lceil \text{cnt} \times 4/3 \rceil$ rounded to the next power
of two (75\% load factor), resizing both up and down between transitions.
This keeps sequential scan cost proportional to state count at all times.

A critical bug was found during this refactoring: \texttt{ht\_clear} zeroed
only the \texttt{keys} array but not \texttt{vals}.
Since \texttt{ht\_add} accumulates with \texttt{+=}, reusing a slot with a
stale value produced incorrect counts.
After fixing this, v2 achieved a DP time of approximately $71$~seconds for
$n = 57$.

% -------------------------------------------------------
\section{Profiling and Multi-Edge Batching: v3}
\label{sec:v3}
% -------------------------------------------------------

\subsection{Per-step profiling}

We added optional per-step instrumentation (elapsed time, state count,
frontier size, and back-edge count) to the C function.
The profile for $n = 57$ revealed:

\begin{center}
\begin{tabular}{lrr}
\toprule
$\nb$ & Steps & Total time\\
\midrule
1 & 5 & $\approx 21$\,s\\
2 & 8 & $\approx 51$\,s\\
3 & 10 & $\approx 283$\,s\\
4 & 1 & $\approx 90$\,s\\
\bottomrule
\end{tabular}
\end{center}

Steps with three or four back-edges at large state counts dominated $70\%$
of runtime.
At $4$~million states and $\nb = 3$, the sequential approach scanned the
source table \emph{three times} — one pass per back-edge.
Each pass reads $\sim\!100$~MB of hash-table data, all DRAM-bound.

\subsection{Multi-edge subset batching}

Instead of $\nb$ sequential table passes, we enumerate all $2^{\nb}$ subsets
of back-edges in a \emph{single} source-table scan.
For each state, each subset is applied sequentially; invalid subsets (degree-2
violations, cycles) are pruned early.
This reduces source-table scans from~$\nb$ to~$1$, giving a $2$--$3\times$
improvement on the bottleneck steps.

The threshold for switching to multi-edge batching is $\nb \ge 2$ and
$\text{states\_in} \ge 200{,}000$.

% -------------------------------------------------------
\section{Fused Edge-and-Elimination Sweeps: v4}
\label{sec:v4}
% -------------------------------------------------------

\subsection{The intermediate table problem}

Even after multi-edge batching, the profiler revealed two catastrophic steps
($156$~s and $201$~s) for $n = 57$.
Investigation showed a two-step cascade:
\begin{enumerate}[nosep]
  \item Step 41 ($\nb = 3$, $3.2$M states): the pre-allocated output table
        (at $1.5\times$ input) is filled to $81\%$ load, triggering an
        internal \texttt{ht\_add} resize to a $9.6$M-slot table.
  \item Step 42 ($\nb = 3$, $3.9$M states): this $9.6$M-slot table becomes
        \texttt{curr}.  Scanning $9.6$M slots to find $3.9$M occupied entries
        is a $2.5\times$ scan waste.  Combined with the elimination that
        follows, total cost was $201$~s.
\end{enumerate}

\subsection{Fused sweeps}

The root cause is the \emph{large intermediate table} between the edge
processing (step B) and the elimination (step C).
The fix is to fuse B and C into a single source-table scan per step:

\begin{enumerate}[nosep]
  \item For each state in \texttt{curr}, enumerate all $2^\nb$ edge subsets;
  \item For each valid result, immediately apply all eliminations due at this
        step via \texttt{apply\_\allowbreak elim\_\allowbreak seq};
  \item Write the final (post-elimination) result directly to \texttt{nxt}.
\end{enumerate}

This eliminates the intermediate table entirely.
The output table is sized for the \emph{final} (post-elimination) state count,
which is typically much smaller than the intermediate.

\textbf{Result.} For $n = 57$, DP time dropped from $73$~s to $42.7$~s
(a $1.7\times$ improvement on the previous best), and the two pathological
steps are no longer present.

% -------------------------------------------------------
\section{Secondary Vertex-Order Refinement}
\label{sec:ordering2}
% -------------------------------------------------------

\subsection{Motivation}

The MaxSAT solver finds a vertex ordering that minimises the maximum frontier
width (i.e.\ the exact pathwidth).
However, many orderings achieve the same pathwidth, and their DP costs differ
dramatically.
The profiles for two distinct $\pw = 11$ orderings of $G_{57}$ showed a
$15\times$ difference in DP time at $n = 35$.

\subsection{The cost function}

The fundamental cost driver per step is:
\begin{equation}
  \text{step cost} \approx 2^{\nb} \times |\text{curr}|
  \label{eq:step_cost}
\end{equation}
where $|\text{curr}|$ is the source table size (which tracks the state count)
and $2^{\nb}$ is the number of edge subsets in the fused sweep.

Two additional effects matter:
\begin{enumerate}[nosep]
  \item \textbf{Table stays large across consecutive expanding steps.}
        When a step has $\text{states\_out} > \text{states\_in}$, the table
        does not shrink, so the \emph{next} step inherits a large table even
        if its own expansion is small.
        A plateau of six consecutive expanding steps at peak frontier width
        caused every step to scan a $\sim\!30$M-slot table to find $\sim\!20$M
        states — a $1.5\times$ overhead per step.
  \item \textbf{Near-peak steps are expensive even when $f < \pw$.}
        A step with $f = \pw - 1$ and $\nb = 4$ can be more expensive than
        one with $f = \pw$ and $\nb = 2$.
\end{enumerate}

We model the per-step cost with a \emph{pressure} variable that accumulates
when consecutive steps are net-expanding ($\nb > n_{\mathrm{elim}}$) and decays
when steps are net-compressing:
\begin{equation}
  \text{cost}(\sigma) = \sum_{\text{steps}} \mathrm{pressure}(s) \times 2^{\nb(s)}
    \times \left(\frac{f(s)}{\pw}\right)^2 + 0.001 \times \text{profile}
  \label{eq:cost_fn}
\end{equation}
where $\text{profile} = \sum_s f(s)$ is the sum of frontier widths (a
classical secondary optimisation criterion).

\subsection{Simulated annealing refinement}
\label{sec:sa}

Given the exact pathwidth ordering from MaxSAT, we run simulated annealing
(random vertex swaps) to minimise the cost function~\eqref{eq:cost_fn} while
rejecting any move that would cause $\max_s f(s) > \pw$.
This preserves the exact pathwidth while optimising the secondary objective.

\begin{table}[H]
\centering
\caption{Effect of SA refinement for $n = 57$, $\pw = 11$.}
\begin{tabular}{lrr}
\toprule
Ordering & DP time & Peak states\\
\midrule
MaxSAT (no refinement) & $42.7$\,s & $4.6$M\\
+ SA (peak-only cost, V1) & $31.8$\,s & $4.6$M\\
+ SA (pressure cost, V2) & $\approx 27$\,s & $4.2$M\\
\bottomrule
\end{tabular}
\end{table}

% -------------------------------------------------------
\subsection{Frontier sojourn and the ordering cost function}
\label{sec:sojourn}
% -------------------------------------------------------

\paragraph{The sojourn problem.}
The cost function~\eqref{eq:cost_fn} penalises the per-step subset
expansion cost $2^{\nb} \times (f/\pw)^2$, but does not explicitly
penalise a vertex that occupies a frontier slot for many consecutive
steps.  We call the number of steps between a vertex's introduction
and its elimination its \emph{frontier sojourn}.

Formally, for a vertex ordering $\sigma$ and vertex $v$, the sojourn is
\[
  \mathrm{soj}(v,\sigma)
  = \max_{w \sim v} \mathrm{pos}_\sigma(w) - \mathrm{pos}_\sigma(v),
\]
where the maximum is taken over the graph-neighbours $w$ of $v$ that
appear later in $\sigma$ than $v$ itself.
The sum $\sum_v \mathrm{soj}(v,\sigma)$ equals the
\emph{profile} $\sum_s f_s(\sigma)$, so the profile term already
present in~\eqref{eq:cost_fn} penalises long sojourns in principle.
The difficulty is one of scale: the profile is
weighted by $0.001$, so a 30-step sojourn at frontier width
$\pw = 14$ contributes $30 \times 14 \times 0.001 = 0.42$ to the
cost — less than $5\%$ of a single peak-width step with
$\nb = 3$ (which contributes $8 \times 1^2 = 8$).
The SA therefore felt no pressure to shorten long sojourns.

\paragraph{Effect on state counts.}
The consequence became visible in the profile for $n = 61$
($\pw = 14$): vertex~$61$ was placed at step~$2$ by the MaxSAT
solver, but its last graph-neighbour appeared at step~$32$.
Vertex~$61$ therefore occupied a frontier slot from step~$2$ through
step~$32$ — a sojourn of $30$ steps, during which the frontier width
was constantly $\pw = 14$.
Had vertex~$61$ been placed at step~$30$ instead, the frontier width
for steps~$2$--$29$ would have been at most $\pw - 1 = 13$.
Since $\Bell(13)/\Bell(14) \approx 14.5\%$, the peak state counts
for those steps would in principle drop from $57$~million to
approximately $8$~million, keeping both hash tables comfortably within RAM.

\paragraph{Huge-page support.}
A secondary cause of slowdown at steps~27--30 of the $n=61$ run
was TLB thrashing.
With both hash tables at ${\sim}128\text{M}$ slots ($3.1$\,GB each),
the working set for random-access hash probing far exceeds the
CPU's TLB (typically $4096\times 4$\,KB pages).
Each probe therefore incurs an additional page-table walk of
${\sim}100$\,ns on top of the DRAM latency, multiplying the observed
cost per state by a factor of $10$--$20$ over the theoretical
prediction.
We address this by calling \texttt{madvise(MADV\_HUGEPAGE)} on every
hash-table allocation of $2$\,MB or larger, reducing the TLB working
set from ${\sim}750{,}000$ entries (at $4$\,KB pages) to
${\sim}1{,}500$ entries (at $2$\,MB pages).

% -------------------------------------------------------
\subsection{Evolution of the cost function}
\label{sec:costfn-evolution}
% -------------------------------------------------------

The SA cost function underwent five iterations before reaching its current
form, each motivated by a specific empirical failure of the preceding version.
Table~\ref{tab:costfn} summarises the progression.

\paragraph{V1 — peak-only penalty.}
The first version penalised only those steps where the frontier width
equals the pathwidth:
\[
  \mathrm{cost}_1(\sigma)
    = \sum_{\{s\,:\,f_s = \pw\}} 2^{\nb(s)}
      + 0.01 \sum_s f_s.
\]
This reduced the $n=57$ DP time from $42.7$\,s to $31.8$\,s.
The failure mode was blindness to near-peak steps: the SA moved an
$\nb=4$ vertex from $f_s = \pw = 11$ (saving $16$ units) into a step
at $f_s = 10$ (costing $16 \times (10/11)^2 = 13.2$ units), which was
operationally more expensive than what was avoided.

\paragraph{V2 — frontier-weighted penalty.}
The second version applied a continuous weight $(f_s/\pw)^2$ to all
steps:
\[
  \mathrm{cost}_2(\sigma)
    = \sum_s 2^{\nb(s)} \cdot \!\left(\tfrac{f_s}{\pw}\right)^{\!2}
      + 0.001 \sum_s f_s.
\]
V2 produced the best result for $n=58$ (138\,s, peak 18.3M states).
Its failure mode emerged at $n=61$: a six-step plateau of consecutive
net-expanding steps at $f_s = \pw$, each scored independently at
$4 \times 1 = 4$ units, was accepted by the SA as equivalent to six
isolated cheap steps, failing to account for the compounding effect on
hash-table size.

\paragraph{V3 — pressure with explicit sojourn penalty.}
The third version added a multiplicative pressure variable and an
explicit per-vertex sojourn penalty (Equation~\ref{eq:cost_fn2_body}
below), intended to address both the plateau problem and long sojourns
simultaneously.
\begin{align}
  \mathrm{pressure}_{s+1} &= \max\!\left(1,\; \mathrm{pressure}_s
      \cdot (1 + 0.3\,\max(0,\, \nb(s) - n_\mathrm{elim}(s)))\right),
      \notag \\
  \mathrm{cost}_3(\sigma)
    &= \sum_s \Bigl[
        \mathrm{pressure}_s \cdot 2^{\nb(s)} \cdot \!\left(\tfrac{f_s}{\pw}\right)^{\!2}
        + \mathrm{soj}(v_s) \cdot \left(\tfrac{f_s}{\pw}\right)^{\!2} \cdot 0.5
      \Bigr]
      + 0.01 \sum_s f_s.
  \label{eq:cost_fn2_body}
\end{align}
The sojourn penalty with weight $0.5$ gave a 30-step sojourn at
$\pw = 14$ a cost of $15$ units — larger than a single peak-width step
with $\nb = 3$.
The SA aggressively moved long-sojourn vertices later, but this
over-correction caused $n=61$ to peak at 423M states instead of 108M:
fixing one vertex's sojourn forced others into positions with four
back-edges at peak width, which was far more expensive.

\paragraph{V4 — pure geometric proxy.}
The fourth version replaced the pressure variable with a purely
multiplicative proxy $\proxy_s$ tracking an estimate of the hash-table
occupancy:
\begin{equation}
  \proxy_{s+1} = \max\!\left(1,\; \proxy_s
      \cdot 2^{\nb(s)} \cdot \!\left(\tfrac{f_s}{\pw}\right)^{\!2}
      \cdot \max\!\left(0.01,\, 0.5^{n_\mathrm{elim}(s)}\right)\right),
  \quad
  \mathrm{cost}_4(\sigma) = \sum_s \proxy_s + 0.001 \sum_s f_s.
\end{equation}
Geometric compounding correctly scores a six-step plateau of $\nb=2$
steps at peak width as $85\times$ the cost of a single such step.
Long sojourns are penalised implicitly: a vertex in the frontier for
$k$ steps participates in $k$ compounding cycles.
The failure mode was that two $\nb=2$ steps are scored identically to
one $\nb=4$ step (since $2^2 \times 2^2 = 2^4$), so V4 could not
distinguish a single expensive step from two cheaper ones.
For $n=58$, this produced a DP time of $265$\,s vs V2's $138$\,s.

\paragraph{V5 — hybrid (current).}
The current version combines both signals by using the proxy as a
\emph{multiplier} rather than as the cost itself:
\begin{equation}
  \proxy_{s+1} = \max\!\left(1,\; \proxy_s
      \cdot 2^{\nb(s)} \cdot \!\left(\tfrac{f_s}{\pw}\right)^{\!2}
      \cdot \max\!\left(0.01,\, 0.5^{n_\mathrm{elim}(s)}\right)\right),
  \label{eq:proxy_update}
\end{equation}
\begin{equation}
  \mathrm{cost}_5(\sigma)
    = \sum_s \proxy_s \cdot 2^{\nb(s)} \cdot \!\left(\tfrac{f_s}{\pw}\right)^{\!2}
      + 0.001 \sum_s f_s.
  \label{eq:cost_fn_v5}
\end{equation}
The step cost $\proxy_s \times 2^{\nb(s)} \times (f_s/\pw)^2$
approximates the actual fused-sweep cost: $\proxy_s$ estimates the
current table occupancy and $2^{\nb(s)}$ is the number of edge subsets
enumerated per source state.
V5 satisfies both desiderata simultaneously:
\begin{itemize}[nosep]
  \item \emph{Single-step accuracy}: $\nb = 4$ is correctly scored
    $2^4/2^2 = 4\times$ more expensive than $\nb=2$ at the same proxy
    level, regardless of history.
  \item \emph{Plateau sensitivity}: four consecutive $\nb=2$ steps at
    $f_s = \pw$ cost $340$ units under V5 vs $4 \times 4 = 16$ under
    V2 — a $21\times$ ratio — correctly reflecting compounding table
    growth.
\end{itemize}
Empirically, V5 produces a $n=59$ run that is more than $2\times$
faster than V2 up to the same point in the computation, with a lower
peak state count (52.6M vs 68.4M at step 36).

\begin{table}[H]
\centering
\caption{Performance of cost function versions across graph sizes.
  DP times are wall-clock seconds; peak states are the maximum hash-table
  occupancy over all steps.
  ``$>2\times$ faster'' for V5 at $n=59$ is relative to V2 and measured
  up to step~36.  Dashes indicate the run was in progress at time of writing.}
\label{tab:costfn}
\small
\setlength{\tabcolsep}{4pt}
\begin{tabular}{llrrrl}
\toprule
Version & Formula summary & $n$ & DP time & Peak states & Note \\
\midrule
No SA & MaxSAT order only
  & 57 & 42.7\,s & 4.6M & Fused B+C baseline \\
V1 & $\sum_{f_s=\pw} 2^{\nb}$
  & 57 & 31.8\,s & 4.6M & $+25\%$ \\
\midrule
No SA & MaxSAT order only
  & 58 & 229\,s  & 22.0M & \\
V1 & $\sum_{f_s=\pw} 2^{\nb}$
  & 58 & 224\,s  & 22.4M & No improvement \\
V2 & $\sum 2^{\nb}(f_s/\pw)^2$
  & 58 & \textbf{138\,s}  & 18.3M & Best for $n=58$ \\
V3 & V2 + pressure + sojourn
  & 58 & 601\,s  & ---   & Regression \\
V4 & $\sum \proxy_s$
  & 58 & 265\,s  & 23.8M & Plateau fix, single-step blind \\
V5 & $\sum \proxy_s \cdot 2^{\nb}(f_s/\pw)^2$
  & 58 & $\approx 186$\,s & --- & Hybrid \\
\midrule
V2 & $\sum 2^{\nb}(f_s/\pw)^2$
  & 59 & --- & 68.4M & Swap at steps 36--37 \\
V5 & $\sum \proxy_s \cdot 2^{\nb}(f_s/\pw)^2$
  & 59 & $>2\times$ faster & 52.6M & Up to step 36 \\
\midrule
V2 & $\sum 2^{\nb}(f_s/\pw)^2$
  & 61 & --- & 108M & Stalls at step 32 \\
V3 & V2 + pressure + sojourn
  & 61 & --- & 423M & $4\times$ worse peak \\
\bottomrule
\end{tabular}
\end{table}

The evolution illustrates a recurring theme: a cost function designed to
capture one phenomenon can inadvertently blind the SA to another.
V5 resolves the tension structurally, by using the geometric proxy as a
multiplier so that both the single-step cost and the accumulated
state-count growth are represented in every term of the sum.

% -------------------------------------------------------
\subsection{Relation to treewidth and the \texttt{htd} library}
\label{sec:treewidth}
% -------------------------------------------------------

A natural question is whether replacing the path decomposition with a
\emph{tree} decomposition of width $\tw(\Gn) \le \pw(\Gn)$ would
reduce the state-space size.
We investigated this using the \texttt{htd} library~\cite{htd},
which computes heuristic tree decompositions via min-fill, min-degree,
and MCS-M elimination strategies with iterative width minimisation.

\paragraph{Treewidth of $\Gn$.}
For $n \le 49$, \texttt{htd}'s min-fill heuristic finds decompositions
of width equal to the (exact) pathwidth, confirming $\tw(\Gn) =
\pw(\Gn)$ in this range.
For $n = 50$ and $n = 55$, it finds width $\pw(\Gn) - 1$
(specifically, $\tw \le 9 < \pw = 10$ and $\tw \le 10 < \pw = 11$),
confirming a gap of at least~$1$.
For $n \ge 57$, the heuristic fails to match the MaxSAT pathwidth
solver and returns decompositions wider than $\pw(\Gn)$; our path
decomposition is strictly better in these cases.

\paragraph{Why the tree-DP approach does not help.}
Even where a gap exists, the improvement is illusory.
The path DP processes each step in $O(\Bell(\pw))$ time.
A tree DP on a width-$\tw$ decomposition introduces \emph{join nodes}
where two independent partial DPs over the same bag must be merged.
For Hamiltonian-path counting, merging requires checking every pair of
compatible connectivity patterns, at cost $O(\Bell(\tw)^2)$ per join
node.
For $\tw = 9$ vs $\pw = 10$:
\begin{align*}
  \text{join-node cost} &= \Bell(9)^2 = 21{,}147^2 \approx 447\text{M}, \\
  \text{path-step cost} &= \Bell(10) = 115{,}975,
\end{align*}
making each join node $3{,}856\times$ more expensive than a path step.
Moreover, $\Gn$ has a near-linear structure (each vertex connects to a
bounded neighbourhood in value), so any tree decomposition resembles a
path decomposition and provides few join nodes to exploit.
The pathwidth-based path DP is optimal for this problem family.

% -------------------------------------------------------
\section{Exact Pathwidth and Optimisation over All Optimal Orderings}
\label{sec:kpw}
% -------------------------------------------------------

The MaxSAT ordering of Section~\ref{sec:ordering} and the annealing of
Section~\ref{sec:ordering2} have since been replaced by a solver built on
Kobayashi's dynamic program for pathwidth~\cite{Kob15}.
It computes the exact pathwidth in seconds, and, more importantly, it
represents a very large family of pathwidth-optimal orderings in a compact
data structure over which a secondary cost can be minimised \emph{exactly}.
The details, with proofs, are in a separate paper~\cite{MillerPathwidth};
here we summarise what the frontier DP uses.

\paragraph{One-pass search.}
Kobayashi's algorithm decides whether $\pw(G)\le k$ by a dynamic program over
\emph{relevant} vertex sets (sets closed under a natural subsumption
expansion), of which there are at most $3^{\tau(G)}$, where $\tau$ is the
vertex cover number~\cite{Kob15,CLTV13}.
The cost of one step of this dynamic program does not depend on~$k$, so the
vertex separation number is the value of a minimum-bottleneck path in a DAG
on relevant sets, and a bucketed Dijkstra-type search finds it in a single
pass, with an optimal ordering, without any bound~$k$ being supplied.
This pass never does more work than one call of the decision procedure at
$k=\pw(G)$.
A greedy rule from the branch and bound of Coudert, Mazauric and
Nisse~\cite{CMN14} can be added to the expansion without affecting any
guarantee; it reduces the number of sets $2$--$6$ times.
On the benchmark graphs of~\cite{MillerPathwidth} the C implementation is
$6$--$8$ times faster than the Sage implementation of that branch and bound,
and $10$--$22$ times faster with the greedy rule.

\paragraph{The DAG of optimal orderings.}
If the search completes the last level instead of stopping at the first
optimal ordering, and then discards the sets from which the full vertex set is
no longer reachable within width $\pw(G)$, the remaining relevant sets form a
DAG every source-to-sink path of which is a pathwidth-optimal ordering.
For $G_{57}$, $G_{58}$, $G_{59}$ and $G_{61}$ the DAG has $13{,}762$,
$45{,}056$, $151{,}041$ and $263{,}672$ nodes and represents at least
$1.6\cdot10^{26}$, $3.6\cdot10^{28}$, $2.0\cdot10^{31}$ and
$2.2\cdot10^{33}$ orderings (each expansion block counted once).
This is a strict subfamily of all pathwidth-optimal orderings: in the frontier
DP's direction, the boundaries between blocks are exactly the prefixes in
which every introduced vertex has an eliminated neighbour.
On small graphs the restriction costs little: for $25\le n\le35$ the optimum
of the V2 cost of Section~\ref{sec:ordering2} over the DAG is within
$0.3$--$1.7\%$ of its optimum over all pathwidth-optimal orderings.

\paragraph{Exact secondary optimisation.}
Every cost used in Section~\ref{sec:ordering2}, and the fitted work model,
is local to one step of the frontier DP: a step contributes $a+p\,e$ to the
cost, where $p$ is the carried state-count proxy, and multiplies $p$ by
$e\,k$.
Such a cost is minimised exactly over all paths of the DAG by a dynamic
program that keeps, at each node, the Pareto-optimal pairs (cost so far,
carried proxy).
This replaces the simulated annealing: the ordering used for each run is the
exact optimum of the chosen cost over the DAG, not a local optimum.
For $G_{75}$ the work-model optimum needs at most $884$ labels at a node.
On graphs with very many equivalent orderings the label sets can grow large
(about $2200$ per node for the $9\times9$ grid), so by default the optimiser
runs with label caps of $64$, $512$, $4096$ and then none, stopping at the
first untruncated (hence exact) result or when a time budget would be
exceeded, and reports whether the result is certified optimal.

\paragraph{Effect.}
For $n=75$ the April pipeline (MaxSAT followed by annealing) took
$2179$\,s to produce its ordering; on the 64-thread Linux machine the new
solver finds the pathwidth in $0.8$\,s and the exact optimum of the work model
over the DAG in $22$\,s.
The DAG also makes the choice among optimal orderings measurable: for
$50\le n\le56$, the total number of DP states over $12$ distinct pathwidth-optimal
orderings per~$n$ (optima of three costs and samples from the DAG) varies
by a factor of $1.4$--$3.0$ without pruning and $1.4$--$2.2$ with it.
Which cost picks the best ordering is not yet settled; at $n=75$ the $V_7$
optimum needed $24\%$ fewer states than the work-model optimum.
Two different optima from the DAG also gave the independent confirmation of
$H(75)$ in Section~\ref{sec:results}.

% -------------------------------------------------------
\section{Memory Efficiency}
\label{sec:memory}
% -------------------------------------------------------

For $n = 59$ ($\pw = 13$), the original v2 implementation approached the
$16$\,GB memory limit.
Two changes addressed this:

\paragraph{Elimination of the \texttt{used\_list}.}
The original hash table stored an auxiliary array of occupied slot indices to
enable O(\textit{cnt}) clearing.
This cost an extra $8$ bytes per slot (matching the key array size).
Replacing it with sequential \texttt{memset} of the keys \emph{and} vals
arrays (necessary since \texttt{ht\_add} uses \texttt{+=}) eliminates the
auxiliary array entirely.
At the load factor change from $50\%$ to $75\%$, the combined effect reduces
memory per state from $32$ to $24$ bytes — a $25\%$ reduction — and at peak
state counts the two-table working set drops from $\sim\!10$\,GB to
$\sim\!4$\,GB.

\paragraph{Right-sizing after fused sweeps.}
When an internal \texttt{ht\_add} resize occurs during a fused sweep, the
output table ends up $2\times$ larger than needed.
Without correction this oversized table becomes \texttt{curr} for the next
step, which then scans it at full (sparse) capacity.
The function \texttt{ht\_compact} detects this and rehashes into a right-sized
table immediately after the sweep, before the table swap.

% -------------------------------------------------------
\section{Extendible Hashing and the Storage Hierarchy}
\label{sec:eh}
% -------------------------------------------------------

The implementation of Sections~\ref{sec:c}--\ref{sec:memory} stores~\nxt{} in a
linear-probing table sized in advance.
At $n=61$ ($\pw=14$, peak $6.2\cdot10^{8}$ states) this fails on a $16$\,GB
laptop: the table must be sized for an estimate of the output, and at the
peak step the old table, the new table and the resize copy are live at once.
Four changes, each addressing a failure observed in the previous version,
made $n=61$ run in $97$ minutes in under $2$\,GB on an Apple M3 Pro.

\paragraph{Extendible hashing.}
We replace the linear-probing table by an extendible hash table
\cite{Fagin}: a directory indexed by the top $gd$ bits of the hash points to
buckets of $B=32$ or $64$ entries, each with a local depth $\ell\le gd$, and
a full bucket is split in two rather than the whole table being rehashed.
Memory is therefore proportional to the actual number of output states, and
there is no resize spike.
The directory entries that point to a bucket of local depth~$\ell$ form an
aligned contiguous range of length $2^{gd-\ell}$, so a split need only
redirect the upper half of that range.
A first implementation scanned the whole directory on each split, which at
$gd=18$ cost about $3.7\cdot10^{10}$ iterations in one step and made the
table ten times slower than linear probing; redirecting only the range makes
the cost of all splits $O(N/B)$.

\paragraph{Radix-sort insertion buffer.}
Buckets are allocated in insertion order, so directory entries with a common
hash prefix point to buckets scattered through memory, and a random insert
costs a DRAM miss.
New pairs are therefore first appended to one of $2^R$ slots selected by the
top $R$ bits of the hash.
When the buffer is full the slots are flushed in order; all entries of one
slot address the same range of the directory, so each bucket is brought into
cache once per flush rather than once per insertion.

\paragraph{Cache-adaptive constants.}
The bucket size~$B$, the state count at which the buffer is used, and $R$
are computed from the L1, L2 and L3 sizes of the machine (read from
\texttt{sysctl} or \texttt{/sys}) and passed to the compiler, so that a bucket
fits in L1 and the region touched by one slot fits in a fraction of L2.

\paragraph{Returning memory to the operating system.}
Even with these changes, a run of $n=61$ reached $70$--$90$\,GB of swap
after about thirty steps.
Memory released with \texttt{free()} returns to the allocator, not to the
operating system, and \texttt{madvise(MADV\_FREE)} is only advisory on macOS.
Freeing each slab of buckets (about $3$\,MB, above the threshold at which
\texttt{malloc} uses \texttt{mmap}) at the end of each step, and reusing one
flat allocation for the radix buffer, causes the pages to be unmapped and
keeps the resident set below $2$\,GB.

% -------------------------------------------------------
\section{Parallel Sort--Merge}
\label{sec:sm}
% -------------------------------------------------------

The values for $62\le n\le 75$ were computed on a Linux machine with
$768$\,GB of memory and $64$ cores, with a parallel variant of the DP step
that replaces the hash table by sorting and merging.
This section describes that variant.

\subsection{Earlier attempts}

Two parallelisations based on hash tables were tried first and abandoned.
In the first, threads scanned disjoint parts of~\curr{} and inserted into
private linear-probing tables, which the main thread then inserted
into~\nxt.
That final merge performs the same $O(|\nxt|)$ random inserts as the serial
sweep, and at $n=45$ two threads were $1.65$ times \emph{slower} than one.
In the second, each worker kept a private extendible hash table and the
merge was itself parallel, but every worker still made random writes into a
table much larger than its cache, and each output state was touched by two
random-access passes.
In both cases the expensive part of a step is not computing the new states
but \emph{reducing} them, that is, adding the counts of equal keys, and a
hash table makes that reduction a sequence of random writes.

\subsection{Independence of contributions}

Let step~$t$ introduce vertex~$v_t$ with $\nb$ back-edges.
Each pair consisting of an entry $(k,c)$ of~\curr{} and a subset~$S$ of the
back-edges yields at most one pair $(k',c)$: apply the edges of~$S$ to~$k$
(rejecting the pair if a vertex would get degree~$3$ or a cycle would close),
then apply the eliminations of step~$t$ (adding $c$ to the total if the
state is complete).
These $|\curr|\cdot 2^{\nb}$ contributions are independent of one another;
the only interaction is that contributions with the same key~$k'$ must be
added.
We therefore separate generation from reduction, and perform the reduction
by sorting, which uses memory bandwidth rather than memory latency and
divides into independent pieces.

\subsection{One step}

\paragraph{Introduce.}
Introducing~$v_t$ writes the code for ``no edges yet'' into a new frontier
slot, i.e.\ it sets bits above all occupied slots.
This is done in place by $P$ threads on disjoint parts of~\curr{}, so no copy
of the table is made.
No entry moves, so the table stays in the order of the hashes of the keys
before the introduce (see below).
The versions used for $62\le n\le 74$ instead re-sorted the table by the hash
of the new keys after every introduce.
That sort is unnecessary: the introduce is injective on keys, so nothing is
combined, and the workers do not need sorted input.
It took $40$--$50\%$ of the DP time in those runs.

\paragraph{Generate.}
Worker $i\in\{0,\ldots,P-1\}$ takes the contiguous slice
$[\,i\cdot\lceil|\curr|/P\rceil,\,(i+1)\cdot\lceil|\curr|/P\rceil\,)$ of~\curr{}
and appends its contributions to a private buffer of
\texttt{CAP} entries of $16$ bytes (a $64$-bit key and a $64$-bit count).
When the buffer is full it is \emph{flushed}: sorted by $h(k')$, where $h$ is
a bijective mixing function of the $64$-bit key, by an eight-pass least
significant digit radix sort on $8$-bit digits; adjacent equal keys are
combined; and the result is stored as a sorted \emph{run}.
The peak memory during generation is
\[
  |\curr| + 2P\cdot\texttt{CAP}\cdot 16\ \text{bytes}
\]
(buffer and radix scratch for each worker), independent of the size of the
output.

\paragraph{Where the runs live.}
Before the workers start, the step estimates the memory needed to merge the
runs in RAM as twice the raw output bound,
$2P\cdot\lceil \lceil|\curr|/P\rceil 2^{\nb}/\texttt{CAP}\rceil\cdot\texttt{CAP}
\cdot 16$ bytes, and compares it with the memory left after the operating
system, \curr{} and the worker buffers.
If it fits, the runs are kept in RAM; otherwise each worker appends its runs
to a backing store on SSD, sized for the worst case.

\paragraph{Merge.}
In the RAM case each worker first merges its own runs.
This merge is split into $256$ independent sub-merges by the top byte
of~$h$, each small enough for its heap and input blocks to stay in L2.
The $P$ resulting streams are then merged into~\nxt{} by the parallel merge
below, or, when memory is tight, by a pairwise tree of merges that frees
each input as soon as it is consumed.
In the SSD case all runs of all workers (up to $694$ streams at $n=75$) are
merged into~\nxt{} in a single parallel pass, each stream read through a
buffer of about (last-level cache size)/(number of streams) entries.

\paragraph{Parallel merge by hash range.}
Given $K$ streams sorted by~$h$ and $T$ threads:
\begin{enumerate}[nosep]
  \item take $64$ evenly spaced values of~$h$ from each stream, sort them, and
        choose $T-1$ splitters that divide the samples into $T$ equal parts;
  \item locate every splitter in every stream by binary search;
  \item bound the output of each of the $T$ hash ranges by the sum of its input
        lengths, and assign each range its own slice of~\nxt;
  \item thread~$r$ merges the $K$ pieces in range~$r$ with a heap,
        combining equal keys, into its slice;
  \item compute prefix sums of the actual output lengths and move the slices
        together.
\end{enumerate}
Since $h$ is a function of the key, equal keys lie in the same range, so the
combining of equal keys is exact and no two threads write the same key.

\subsection{Load balancing}

The amount of work generated by a state is the number of subsets~$S$ that
survive the degree, cycle and elimination tests, and this varies by orders of
magnitude between states; nothing cheap to compute from the key predicts it
well.
We did not try to predict it.
Every table produced by the merge is ordered by $h(k)$.
The in-place introduce at the next step changes each key but moves no entry,
so the table is then ordered by the hash of the previous key, which is just
as unrelated to the structure of the new state.
A contiguous slice of~\curr{} is therefore in effect a random sample of the
states.
Static slices are therefore balanced in expectation.
In the $n=75$ run with $64$ workers, the ratio of the slowest worker's
generation time to the mean was between $1.03$ and $1.19$ on all steps that
contribute significantly to the running time; it reached $1.5$--$1.9$ only
on short steps.
A shared queue of small chunks for work stealing was implemented but was
never needed, and is not used.

The same ordering balances the merge.
An earlier version sorted runs by the raw key but chose splitters in hash
space.
For frontier width at most $14$ the top byte of a raw key contains only the
two flag bits, so the runs were
concentrated in $2$ of $256$ buckets, the binary searches for hash
splitters in a stream sorted by raw key returned meaningless positions, and
five of six merge threads did no work.
The bucketed per-worker merge also found wrong bucket boundaries in that
version; that it still returned correct counts was a matter of luck.
Sorting, bucketing, splitting and searching now all use $h$.

\subsection{Safeguards}

\begin{itemize}[nosep]
  \item \emph{Counts.} Reducing the count in an entry from $128$ to $64$ bits
        reduced entries from $24$ to $16$ bytes. Every addition of two counts
        tests for unsigned wraparound and sets a global flag, and the driver
        raises an error if the flag is set. The final total is accumulated in
        $128$ bits.
  \item \emph{Backing store.} An earlier version marked pages of spilled
        runs with \texttt{MADV\_FREE}. Under memory pressure the kernel
        reclaimed them and later reads returned zeros; a zero key sorts below
        every valid bucket and was skipped, so data were lost silently. Live
        run data are no longer marked, and the backing store is sized for the
        worst case.
  \item \emph{False sharing.} The per-worker state is padded to a $128$-byte
        cache line; without this, updates of buffer lengths by different
        workers slowed steps with few runs by a factor of $3$--$4$.
  \item \emph{Checkpointing.} The sort--merge variant does not support
        checkpointing, so a run that fails must be restarted.
\end{itemize}

\subsection{Tuning}

The number of threads, \texttt{CAP}, the sizes of the merge block and stream
buffers, the threshold below which merges are serial, and the last-level
cache size are read from a per-machine configuration file.
On the $768$\,GB machine, the April run of $n=75$ completed only after these
parameters had been adjusted by hand: of fourteen attempts, the others aborted
when memory or storage was exhausted. (Its count was later withdrawn; see
Section~\ref{sec:results}.)
For the orderings, the simulated-annealing cost of
Section~\ref{sec:ordering2} was extended to charge a step that spills to SSD
two to three times as much as a step of the same size that does not.

On the laptop with the same ordering, the sort--merge variant with six
threads computed $n=59$ in $317$\,s, against $385$\,s for the best
extendible-hashing run.
On the Linux machine, $n=61$ takes $699$\,s of DP time, against $97$ minutes
for extendible hashing on the laptop.

% -------------------------------------------------------
\section{Pruning States That Cannot Be Completed}
\label{sec:prune}
% -------------------------------------------------------

The idea of this section is taken from a program of Zhao Hui
Du~\cite{DuPath4} that counts the Hamiltonian paths of~$\Gn$.
His program processes the vertices in the fixed order $n,n-1,\ldots,1$ and
discards a partial state as soon as the vertices that are not yet finished
can no longer be joined, together with the existing path fragments, into a
single path.
It does this with a degree test and a connectivity test.
Our DP had no such test.
A state that can never be completed is carried, multiplied by up to
$2^{\nb}$ at each step, sorted and merged, until it is rejected at an
elimination or at the last step.
The simplest example is a state with three half-external chains
(Section~\ref{sec:dp}): no two of them can ever be joined before the last
step, and no single path has three ends outside the frontier.

\subsection{Necessary conditions}

After step~$t<n-1$ let $F$ be the frontier, after the eliminations of
step~$t$, and $U=\{v_{t+1},\ldots,v_{n-1}\}\neq\emptyset$ the vertices not yet
introduced.
Every edge with no end in~$U$ has been decided.
For a state $\sigma$ let $N=\{u\in F:\sigma_u\neq-1\}$ be the frontier
vertices that are not interior, and for $u\in F$ let
$d_U(u)=|\Gamma(u)\cap U|$, where $\Gamma(u)$ is the set of neighbours
of~$u$; this is at least~$1$ since otherwise $u$
would have been eliminated.
Let $H_\sigma$ be the graph on $N\cup U$ whose edges are the edges of~$G$
inside~$U$, the edges of~$G$ between $N$ and~$U$, and one \emph{link}
joining the two frontier ends of each chain that has both ends in~$F$.

\begin{proposition}\label{prop:prune}
If a state $\sigma$ after step $t<n-1$ extends to a Hamiltonian path of~$G$,
then
\begin{enumerate}[nosep,label=(\Alph*)]
  \setcounter{enumi}{2}
  \item $N\neq\emptyset$ and $H_\sigma$ is connected;
  \item every $w\in U$ has $\deg_{H_\sigma}(w)\ge1$, and
  \[
    e(\sigma) = m(\sigma)
      + \#\{u\in F:\sigma_u=0,\ d_U(u)=1\}
      + \#\{w\in U:\deg_{H_\sigma}(w)=1\}\ \le\ 2,
  \]
  where $m(\sigma)$ is the number of half-external chains.
\end{enumerate}
\end{proposition}

\begin{proof}
Let $P$ be a Hamiltonian path completing~$\sigma$.
Its edges are the decided edges of~$\sigma$ and a set $E'$ of edges with an
end in~$U$; an edge of~$E'$ cannot meet an interior or eliminated vertex,
so $E'\subseteq E(H_\sigma)$.
The decided edges form vertex-disjoint paths (the chains, and the trivial
paths at vertices with $\sigma_u=0$).
No chain has both ends eliminated, since then it would be all of~$P$
and $U$ would be empty.
So each chain has either both ends in~$N$, or one end in~$N$ and the other
an eliminated end of~$P$.
Replacing each chain of the first kind by its link turns $P$ into a
spanning path of $N\cup U$ in $H_\sigma$, which proves~(C); $N$ is non-empty
because $P$ contains an edge between a decided vertex and~$U$.
For~(D): every vertex of~$P$ has degree at least~$1$ in~$P$.
A vertex $w\in U$ has $\deg_P(w)\le\deg_{H_\sigma}(w)$, and a vertex
$u\in F$ with $\sigma_u=0$ has $\deg_P(u)\le d_U(u)$, so the vertices in the
last two terms of $e(\sigma)$ are ends of~$P$; so is the eliminated end of
each half-external chain.
These vertices are distinct, and $P$ has two ends.
\end{proof}

A Hamiltonian cycle found at the last step minus its closing edge, which
is decided at that step, is a Hamiltonian path completing the same state,
so the proposition also justifies the pruning when cycles are counted.
Dropping every state that violates (C) or (D) therefore leaves both counts
unchanged.

\subsection{Implementation}

At each step the components of $G[U]$ are computed once by union--find,
and for each frontier slot $u$ we record $d_U(u)$ and the bit set of the
components containing a neighbour of~$u$; for each $w\in U$ with a frontier
neighbour we record $\deg_{G[U]}(w)$ and the bit set of its frontier
neighbours.
For a key, $e(\sigma)$ is then a few population counts, and~(C) is a
fix-point iteration on two bit sets (slots reached and components reached),
in time $O(\fs\cdot\text{iterations}+|\partial U|)$.
The test is applied by the workers of Section~\ref{sec:sm} to every
contribution after the eliminations and before it enters a buffer, so a
dead state is never sorted, never merged and never generates successors.
Du's program also commits an edge as soon as an endpoint has a single
possible neighbour left; that has no counterpart here, where an edge is
decided when its later end is introduced.

\subsection{Effect}

The counts with pruning agree with those without it, and with
A071983~\cite{OEIS-A071983}, for every $n$ with $15\le n\le 50$ and for
$n=53,55,56,57,58$; the cycle counts agree for $32\le n\le 48$; and the
counts agree with Held--Karp on random graphs of in-step width~$16$ and on
graphs with many back-edges.
At $n=40$ between $25\%$ and $75\%$ of the states of each step are dead,
and at $n=52$, summed over all steps, $54.6\%$ are.
Every rejected state violates~(D); on these graphs and orderings
condition~(C) never rejected a state that (D) accepted.
Table~\ref{tab:prune} compares the DP times with and without pruning on a
four-thread Linux virtual machine; at $n=57$ pruning halves the largest
table, and at step~$37$ of $n=58$ it reduces the table from $24.5$ to $7.5$
million states.
The cleanest comparison is Table~\ref{tab:prune61}: $n=61$ on an M3~Pro
laptop with $18$\,GB, twice with the same build, configuration and
ordering, once without and once with pruning.
Pruning reduces the DP time by a factor of $4.2$, and the states summed
over all steps by $4.4$.
Generating the contributions gains least ($3.4\times$), since every key is
tested and the dropped contributions are still generated; sorting and
merging gain more than the number of states ($7.4\times$ and $8.7\times$),
because without pruning eleven steps did not fit in RAM and wrote their
runs to SSD.
Peak memory falls by only $1.4\times$: the worker buffers
($2P\cdot\texttt{CAP}$ entries, $6.4$\,GB here) are the same in both runs.
The gain grows with~$n$, from $1.5$--$2.3\times$ at $55\le n\le 57$ to
$4.2\times$ at $n=61$.
The ratio of table sizes at a step is not a product over earlier steps:
the unpruned table holds every reachable state, including all successors
of states that are already dead, and the pruned table holds the states
that pass the tests at every step so far.
That ratio grows along the run (at $n=61$ it is $4.9$ at the largest
step) because the proportion of dead states grows as vertices are
eliminated and path ends accumulate outside the frontier.
All of these are single runs.
% Further timings (n=58 on the iMac, larger n) to be added.
The runs for $62\le n\le 75$ reported in Section~\ref{sec:results} were made
without pruning.

\begin{table}[H]
\centering
\caption{DP time with and without pruning (Linux VM, $4$ threads, one run
each, same ordering in each row).}
\label{tab:prune}
\begin{tabular}{rrrr}
\toprule
$n$ & without & with & ratio\\
\midrule
55 & $9.4$\,s & $5.2$\,s & $1.8$\\
56 & $16.6$\,s & $11.2$\,s & $1.5$\\
57 & $68.2$\,s & $29.4$\,s & $2.3$\\
58 & $>174$\,s & $130.1$\,s & --\\
\bottomrule
\end{tabular}
\end{table}
\noindent
For $n=58$ the run without pruning was stopped after step~$37$, at $158.0$\,s;
the run with pruning reached that step at $64.6$\,s.

\begin{table}[H]
\centering
\caption{$n=61$ with and without pruning (M3 Pro, $18$\,GB, $6$ threads,
same build and ordering, pathwidth $14$; one run each).
Phase times are sums over the steps of the slowest worker.}
\label{tab:prune61}
\begin{tabular}{lrrr}
\toprule
 & without & with & ratio\\
\midrule
DP time & $1177.1$\,s & $282.0$\,s & $4.17$\\
contributions generated & $15.95$\,G & $4.39$\,G & $3.6$\\
\quad of which dropped by pruning & -- & $1.88$\,G & \\
states, summed over steps & $6.52$\,G & $1.50$\,G & $4.4$\\
largest table (step 41) & $702$\,M & $144$\,M & $4.9$\\
peak resident memory & $12.08$\,GB & $8.66$\,GB & $1.4$\\
steps spilling runs to SSD & $11$ & $0$ & \\
generate & $773$\,s & $228$\,s & $3.4$\\
sort (worker flushes) & $236$\,s & $32$\,s & $7.4$\\
merge & $182$\,s & $21$\,s & $8.7$\\
\bottomrule
\end{tabular}
\end{table}

% -------------------------------------------------------
\section{Results}
\label{sec:results}
% -------------------------------------------------------

\subsection{Computation}

Table~\ref{tab:results} summarises the evolution of the per-step timing
improvements.

\begin{table}[H]
\centering
\caption{DP time for $n = 57$ across implementation versions.
All timings use the same exact MaxSAT pathwidth ordering ($\pw = 11$).}
\label{tab:results}
\begin{tabular}{llr}
\toprule
Version & Key change & DP time ($n=57$)\\
\midrule
v1 & C with byte-array keys, malloc/free per sweep & $\approx 134$\,s\\
v2 & Packed \texttt{uint64\_t} keys, dynamic table sizing & $71$\,s\\
v3 & Multi-edge subset batching & $73$\,s$^*$\\
v4 & Fused edge-and-elimination sweep & $42.7$\,s\\
v4 + SA V1 & Peak-only cost function & $31.8$\,s\\
v4 + SA V2 & Frontier-weighted cost function & $\approx 27$\,s\\
v4 + SA V5 & Hybrid geometric proxy (current) & $\approx 25$\,s\\
\bottomrule
\multicolumn{3}{l}{$^*$v3 alone added overhead at small steps; the gain appears at $n \ge 55$.}
\end{tabular}
\end{table}

\subsection{Hamiltonian path counts}

Table~\ref{tab:counts} lists the counts obtained, which form OEIS
sequence A071983.
Values for $n\le 61$ were computed with the extendible-hashing version on a
laptop; the values for $57\le n\le 61$ were reproduced by the sort--merge
version on the Linux machine.
Values for $62\le n\le 75$ were computed only with the sort--merge version.
All DP times in the table are for the sort--merge version on the Linux
machine, and exclude the ordering computation
(MaxSAT followed by simulated annealing, $5$--$37$ minutes).
The values for $62\le n\le 74$ were computed with the code of mid-April,
which used $24$-byte entries and re-sorted the table after each introduce;
the value for $n=75$ was computed with the current code ($16$-slot keys and
pruning, Section~\ref{sec:prune}), twice, with two different optimal orderings
(see below).
For $n=74$ four runs gave the same count: the first with the earlier code
($3449$\,s, shown in the table), and three with the later code ($945$, $957$
and $1205$\,s).
All four used the same ordering, so they check the code changes but not the
ordering.
Of the $2500$\,s saved, about $1400$\,s is the removed sort after the
introduce, $760$\,s is merging (including three steps that no longer spill to
SSD because entries are smaller), and $280$\,s is flushing; the time spent
generating states did not change.
The times for $62\le n\le 73$ would presumably fall by a similar factor with
the later code.

\begin{table}[H]
\centering
\caption{Undirected Hamiltonian paths in~$G_n$. $\pw$ is the width of the
ordering used, which is the pathwidth.}
\label{tab:counts}
\begin{tabular}{rrrr}
\toprule
$n$ & $\pw$ & $|\text{Ham}(G_n)|$ & DP time\\
\midrule
56 & & 47{,}746{,}712{,}739 & \\
57 & 11 & 68{,}092{,}497{,}615 & 27\,s\\
58 & 12 & 123{,}092{,}214{,}818 & 86\,s\\
59 & 13 & 332{,}409{,}010{,}079 & 73\,s\\
60 & 13 & 783{,}175{,}077{,}980 & 202\,s\\
61 & 14 & 3{,}711{,}439{,}128{,}718 & 699\,s\\
62 & 14 & 19{,}133{,}360{,}717{,}991 & 1049\,s\\
63 & 14 & 121{,}460{,}168{,}051{,}419 & 1162\,s\\
64 & 14 & 329{,}054{,}896{,}751{,}390 & 1551\,s\\
65 & 14 & 780{,}444{,}607{,}816{,}976 & 1643\,s\\
66 & 14 & 1{,}389{,}466{,}102{,}344{,}323 & 1858\,s\\
67 & 14 & 2{,}317{,}807{,}210{,}864{,}530 & 1742\,s\\
68 & 14 & 5{,}019{,}822{,}900{,}252{,}036 & 2167\,s\\
69 & 14 & 6{,}697{,}388{,}275{,}608{,}246 & 1465\,s\\
70 & 14 & 12{,}684{,}480{,}251{,}402{,}254 & 1635\,s\\
71 & 14 & 65{,}830{,}124{,}660{,}217{,}873 & 2795\,s\\
72 & 14 & 132{,}407{,}660{,}925{,}860{,}729 & 2418\,s\\
73 & 14 & 536{,}884{,}024{,}745{,}221{,}467 & 2079\,s\\
74 & 14 & 2{,}176{,}025{,}661{,}130{,}535{,}741 & 3449\,s\\
75 & 15 & 9{,}231{,}885{,}593{,}047{,}053{,}408$^{\dagger}$ & 1047\,s\\
\bottomrule
\multicolumn{4}{l}{$^{\dagger}$Current code, with pruning; confirmed by a second ordering (see below).}
\end{tabular}
\end{table}

\paragraph{The value for $n=75$.}
$n=75$ is the first case in which the frontier reaches width~$15$.
During a step the frontier briefly holds one more vertex, the one being
introduced, so a width-$15$ ordering needs $16$ slots.
The key encoding used in April had fifteen $4$-bit slots and two flag bits;
the sixteenth slot overlapped the flag bits, and the count it produced,
$652{,}211{,}775{,}281{,}872{,}067{,}084$, is wrong and is withdrawn.
It was out of line with its neighbours ($H(75)/H(74)\approx300$, whereas
$H(n)/H(n-1)$ lies between $1.3$ and $6.4$ for $62\le n\le74$).
The present encoding uses sixteen $4$-bit slots and no flag bits; the flags
are not needed (a closed chain before the last step is rejected at once, and
a closure at the last step is counted at once).
It agrees with Held--Karp on random graphs whose orderings need exactly
$16$ slots.
With it, two runs of $n=75$ with different pathwidth-optimal orderings
(the work-model optimum and the $V_7$ optimum, which differ in $71$ of the
$75$ positions) both give $H(75)=9{,}231{,}885{,}593{,}047{,}053{,}408$,
in $1047$ and $963$\,s on the $64$-core machine.
Now $H(75)/H(74)\approx4.24$, and the value lies on the asymptotic predictions
of the companion paper (residual $-0.34$ against the constant-$C$ formula and
$+1.99$ against the scaling predictor, against $-0.75$ and $+1.95$ at $n=74$).
The count is below $2^{64}$; no exact count has yet exceeded $64$ bits, although
intermediate counts did before pruning (Section~\ref{sec:prune}), and their
wraps were recovered exactly by the Chinese remainder theorem.

\subsection{Alternative formula}

We also evaluated the inclusion-exclusion formula of Giscard and
Rochet~\cite{GiscardRochet} for counting Hamiltonian paths via connected
dominating sets.
Empirically, the number of connected dominating sets of~$\Gn$ grows as
approximately $3^n$, making the formula impractical for $n \ge 30$ on
$\Gn$ (estimated $4 \times 10^{20}$ terms for $n = 57$).
The formula is effective for graphs with small dominating sets (e.g.\ social
networks with hub vertices) but not for graphs like~$\Gn$ where the minimum
dominating set has size $\approx 5n/6$.

% -------------------------------------------------------
\section{Conclusion}
\label{sec:conclusion}
% -------------------------------------------------------

We have described an iterative process of algorithmic engineering that extended
the frontier of exact Hamiltonian path counting in square-sum graphs from
$n = 56$ (limited by ZDD memory) to $n = 75$.

The main technical contributions are:
\begin{enumerate}[nosep]
  \item A clean formulation of the frontier DP with a correctness fix for the
        closed-chain counting bug.
  \item A compact C implementation using packed \texttt{uint64\_t} state keys
        and dynamic table sizing with right-sizing after sweeps.
  \item Fused edge-and-elimination sweeps that eliminate the large intermediate
        table that was the dominant runtime bottleneck.
  \item A secondary vertex-ordering refinement via simulated annealing with a
        hybrid geometric-proxy cost function (V5, Equation~\ref{eq:cost_fn_v5})
        that correctly models both single-step table-scan cost and cumulative
        state-count growth across consecutive expanding steps.
  \item TLB-aware memory allocation via \texttt{madvise(MADV\_HUGEPAGE)},
        eliminating the page-table bottleneck at peak state counts.
  \item Checkpoint/restart with an ordering hash guard, preventing silent
        corruption when a checkpoint from one vertex ordering is inadvertently
        loaded by a run with a different ordering.
  \item An extendible hash table with a radix-sorted insertion buffer and
        cache-adaptive constants, whose memory is proportional to the actual
        output, and explicit return of freed slabs to the operating system.
  \item A parallel sort--merge DP step: independent contributions are
        generated by workers on static slices of the table, reduced by radix
        sorting into bounded runs held in RAM or on SSD, and merged by threads
        that own disjoint ranges of hash space.
        Keeping every table in hash order balances the work without a cost
        model.
  \item An ordering solver based on Kobayashi's dynamic program: exact
        pathwidth in one pass, a DAG representing about $10^{26}$--$10^{33}$
        pathwidth-optimal orderings of~$\Gn$, and exact minimisation of a
        step-local secondary cost over that DAG~\cite{MillerPathwidth}.
  \item Pruning, after Zhao Hui Du, of the states that can no longer be
        completed to a Hamiltonian path (an endpoint budget and a
        connectivity test), which reduces the DP time at $n=61$ by a
        factor of $4.2$.
\end{enumerate}

The approach is general and can be applied to any sparse graph of moderate
pathwidth: the ordering solver of Section~\ref{sec:kpw} computes the exact
pathwidth and an optimised optimal ordering directly.

\paragraph{Acknowledgements.}
The author thanks Claude (Anthropic) for assistance with algorithm design,
implementation, profiling, and debugging throughout this work.

% -------------------------------------------------------
% References
% -------------------------------------------------------
\nocite{Minato1993,Kinnersley,Bellman,HeldKarp}
\bibliographystyle{amsplain}
\bibliography{square_sum}

\end{document}
